\documentclass[../numerical_methods.tex]{subfiles}
\begin{document}
\section{Aula 08 - 26 de Março, 2026}
\subsection{Motivações}
\begin{itemize}
	\item Aplicação: redes sociais.
\end{itemize}
\subsection{Mineração de Dados em Redes Sociais}
Imagine que você foi contratado como Engenheiro de Dados por uma nova rede social corporativa. A diretoria precisa que você construa um sistema modular em Python capaz de ler a matriz de conexões dos usuários para extrair inteligência de negócios dela, incluindo a descoberta dos maiores influenciadores, quem está isolada e criar um algoritmo de ``Pessoas que você talvez conheça''.
Qual a base teórica matemática para fundamentar esse pedido?

Vamos fazer uma mistura de áreas: utilizaremos a álgebra linear e a teoria de grafos; dando um contexto para esta segunda, uma rede de conexões na teoria dos grafos pode ser representada por uma matriz quadrada A de ordem n, junto a algumas regras:
\begin{itemize}
	\item Se o usuário da linha i segue o usuário da coluna j, então \(a_{ij}=1\); caso contrário, \(a_{ij}=0\);
	\item A diagonal principal é sempre 0, afinal seguir a si próprio não conta;
	\item A soma da linha i é exatamente quantas pessoas o usuário i segue; e
	\item A soma da coluna j é exatamente quantas pessoas seguem o usuário j.
\end{itemize}

Para recomendar amigos, precisamos achar amigos em comum, e é aqui que entra a álgebra linear -- o produto da matriz de redes sociais por ela mesma resulta na matriz de conexões indiretas \(C = A \times A\)! O elemento \(c_{ij}\) indica exatamente quantos amigos em comum os usuários i e j possuem.
Àqueles que forem fazer pelo Colab, seguem algumas instruções:
\begin{itemize}
	\item[1)] Preencha a lógica matemática onde está escrito TODO nos códigos;
	\item[2)] O módulo 1, nomeado \texttt{motor\_matrizes}, não pode conter \texttt{prints} ou bibliotecas prontas, apenas matemática pura!;
	\item[3)] Sempre que alterar o código em um módulo, rode a célula novamente para salvar o arquivo no Colab.
\end{itemize}
Após realizar a escrita dos códigos, há um que juntará os 3 módulos e testar se está tudo certo; nele, não altere nada! É um banco de dados fictício para avaliar se deu certo.

\textbf{\underline{Dica de Ouro do Colab}:} se você rodar e der erro, arrume o código na célula de cima, rode-a novamente para salvar o arquivo, clique em \texttt{Ambiente de Execução \> Reiniciar sessão} no menu superior e tente rodar o app novamente!

\begin{listing}[H]
	\caption{\textbf{Redes Sociais:} Motor de Matrizes}
	\begin{minted}[bgcolor = DarkBackground, linenos=true, breaklines]{Python}
%%writefile motor_matrizes.py
# ================================================
# MÓDULO 1: MOTOR DE MATRIZES (A Matemática Pura)
# REGRAS: SEM prints, SEM bibliotecas externas.
# ================================================

def multiplicar_matrizes(A, B):
    """Multiplica duas matrizes quadradas A e B. Retorna a matriz C."""
    # TODO: Implemente a lógica de multiplicação usando laços 'for' aninhados.
    # Dica: Lembre-se do aquecimento de criptografia! Crie uma matriz cheia de zeros primeiro.
    return []

def somar_coluna(M, indice_coluna):
    """Retorna a soma de todos os elementos de uma coluna específica."""
    soma = 0
    # TODO: Crie um laço para somar os valores da coluna solicitada.
    return soma

def somar_linha(M, indice_linha):
    """Retorna a soma de todos os elementos de uma linha específica."""
    soma = 0
    # TODO: Crie um laço para somar os valores da linha solicitada.
    return soma
 \end{minted}
\end{listing}

\begin{listing}[H]
	\caption{\textbf{Redes Sociais:} Algoritmo da Rede Social}
	\begin{minted}[bgcolor = DarkBackground, linenos=true, breaklines]{Python}
%%writefile algoritmo_rede.py
# ===========================================================
# MÓDULO 2: ALGORITMO DA REDE SOCIAL (A Lógica de Negócios)
# ===========================================================
import motor_matrizes as motor

def encontrar_influenciador(matriz_rede):
    """Retorna o ÍNDICE da coluna com a maior soma (mais seguidores)."""
    # TODO: Use um laço e a função 'motor.somar_coluna' para descobrir o maior valor.
    return -1

def encontrar_isolados(matriz_rede):
    """Retorna uma lista com os índices de quem não segue ninguém e não é seguido."""
    isolados = []
    # TODO: Use as funções do motor. Se a soma da linha E da coluna for 0, adicione à lista.
    return isolados

def gerar_recomendacoes(matriz_rede):
    """Retorna a matriz resultante da operação A x A (amigos em comum)."""
    # TODO: Chame a função do motor para multiplicar a matriz por ela mesma.
    return []
 \end{minted}
\end{listing}

\begin{listing}[H]
	\caption{\textbf{Redes Sociais:} Interface de Usuário}
	\begin{minted}[bgcolor = DarkBackground, linenos=true, breaklines]{Python}
%%writefile interface_usuario.py
# =================================================
# MÓDULO 3: INTERFACE DE USUÁRIO (A Visualização)
# =================================================

def exibir_perfil_usuario(nomes, indice_usuario, num_seguidores):
    """Imprime os dados do influenciador da rede."""
    # TODO: Formate e imprima o nome do usuário e a quantidade de seguidores.
    pass

def exibir_isolados(nomes, indices_isolados):
    """Imprime os nomes dos usuários que estão isolados."""
    # TODO: Percorra a lista de índices e imprima os nomes correspondentes.
    pass

def exibir_recomendacoes(nomes, matriz_A2, indice_alvo):
    """Lê a linha do alvo na matriz_A2 e sugere amigos em comum."""
    nome_alvo = nomes[indice_alvo]
    print(f"--- SUGESTÕES DE AMIZADE PARA: {nome_alvo.upper()} ---")

    # TODO: Verifique a linha da matriz_A2 correspondente ao 'indice_alvo'.
    # Imprima quem tem conexões em comum (valor > 0) e que não seja o próprio alvo.
    pass
 \end{minted}
\end{listing}

\begin{listing}[H]
	\caption{\textbf{Redes Sociais:} Painel de Testes}
	\begin{minted}[bgcolor = DarkBackground, linenos=true, breaklines]{Python}
# =============================================
# SCRIPT PRINCIPAL: APP SOCIAL (Test Harness)
# =============================================
import motor_matrizes as math_engine
import algoritmo_rede as rede
import interface_usuario as ui

# 1. Base de Dados
nomes = ["Ana", "Bruno", "Carlos", "Diana", "Eduardo"]

# Matriz de Conexões (A)
matriz_A = [
    [0, 1, 1, 0, 0],  # Ana segue Bruno e Carlos
    [1, 0, 1, 1, 0],  # Bruno segue Ana, Carlos e Diana
    [1, 0, 0, 1, 0],  # Carlos segue Ana e Diana
    [1, 0, 0, 0, 0],  # Diana segue Ana
    [0, 0, 0, 0, 0]   # Eduardo é isolado
]

print("=" * 40)
print(" INICIANDO ANÁLISE DA REDE SOCIAL ")
print("=" * 40 + "\n")

try:
    # 2. Processamento (Lógica)
    indice_famoso = rede.encontrar_influenciador(matriz_A)

    # Proteção: só calcula seguidores se o aluno já programou a função e não retornou -1
    if indice_famoso != -1:
        seguidores_famoso = math_engine.somar_coluna(matriz_A, indice_famoso)
    else:
        seguidores_famoso = 0

    isolados = rede.encontrar_isolados(matriz_A)
    matriz_conexoes = rede.gerar_recomendacoes(matriz_A)

    # 3. Exibição (Interface)
    ui.exibir_perfil_usuario(nomes, indice_famoso, seguidores_famoso)
    print()
    ui.exibir_isolados(nomes, isolados)
    print()

    # Testando recomendações para a Ana (índice 0)
    if len(matriz_conexoes) > 0:
        ui.exibir_recomendacoes(nomes, matriz_A, matriz_conexoes, 0)
    else:
        print("A matriz de recomendações ainda está vazia. Implemente a multiplicação!")

except Exception as e:
    print(f"ERRO NA EXECUÇÃO: {e}")
    print("Verifique se você implementou as funções corretamente nos módulos!")
 \end{minted}
\end{listing}

\subsection{Gabarito das Redes Soicais}
\begin{listing}[H]
	\begin{minted}[bgcolor = DarkBackground, linenos=true, breaklines]{Python}
%%writefile motor_matrizes.py
# ==============================
# GABARITO: MOTOR DE MATRIZES
# ==============================

def multiplicar_matrizes(A, B):
    """Multiplica duas matrizes quadradas A e B. Retorna a matriz C."""
    linhas_A = len(A)
    colunas_B = len(B[0])

    # 1. Criar a matriz C cheia de zeros
    C = []
    for i in range(linhas_A):
        C.append([0] * colunas_B)

    # 2. Multiplicação com 3 laços 'for' aninhados
    for i in range(linhas_A):
        for j in range(colunas_B):
            soma = 0
            for k in range(len(B)):
                soma += A[i][k] * B[k][j]
            C[i][j] = soma

    return C

def somar_coluna(M, indice_coluna):
    """Retorna a soma de todos os elementos de uma coluna específica."""
    soma = 0
    linhas = len(M)
    for i in range(linhas):
        soma += M[i][indice_coluna]
    return soma
\end{minted}
\end{listing}

\begin{minted}[bgcolor = DarkBackground, linenos=true, breaklines]{Python}
def somar_linha(M, indice_linha):
    """Retorna a soma de todos os elementos de uma linha específica."""
    soma = 0
    colunas = len(M[0])
    for j in range(colunas):
        soma += M[indice_linha][j]
    return soma
\end{minted}

\begin{listing}[H]
	\begin{minted}[bgcolor = DarkBackground, linenos=true, breaklines]{Python}
%%writefile algoritmo_rede.py
# ====================================
# GABARITO: ALGORITMO DA REDE SOCIAL
# ====================================
import motor_matrizes as motor

def encontrar_influenciador(matriz_rede):
    """Retorna o ÍNDICE da coluna com a maior soma (mais seguidores)."""
    maior_soma = -1
    indice_famoso = -1
    colunas = len(matriz_rede[0])

    for j in range(colunas):
        soma_atual = motor.somar_coluna(matriz_rede, j)
        if soma_atual > maior_soma:
            maior_soma = soma_atual
            indice_famoso = j

    return indice_famoso

def encontrar_isolados(matriz_rede):
    """Retorna uma lista com os índices de quem não segue ninguém e não é seguido."""
    isolados = []
    ordem_matriz = len(matriz_rede)

    for i in range(ordem_matriz):
        quantos_segue = motor.somar_linha(matriz_rede, i) 
        quantos_seguidores = motor.somar_coluna(matriz_rede, i) 

        if quantos_segue == 0 and quantos_seguidores == 0:
            isolados.append(i)

    return isolados
\end{minted}
\end{listing}

\pagebreak
\begin{minted}[bgcolor = DarkBackground, linenos=true, breaklines]{Python}
def gerar_recomendacoes(matriz_rede):
    """Retorna a matriz resultante da operação A x A (amigos em comum)."""
    # Apenas uma linha de código chama o motor pesado!
    return motor.multiplicar_matrizes(matriz_rede, matriz_rede)
\end{minted}

\begin{listing}[H]
	\begin{minted}[bgcolor = DarkBackground, linenos=true, breaklines]{Python}
%%writefile interface_usuario.py
# ================================
# GABARITO: INTERFACE DE USUÁRIO
# ================================

def exibir_perfil_usuario(nomes, indice_usuario, num_seguidores):
    """Imprime os dados do influenciador da rede."""
    nome = nomes[indice_usuario]
    print(f"MAIOR INFLUENCIADOR: {nome}")
    print(f"Seguidores: {num_seguidores}")

def exibir_isolados(nomes, indices_isolados):
    """Imprime os nomes dos usuários que estão isolados."""
    print("USUÁRIOS ISOLADOS:")
    if len(indices_isolados) == 0:
        print("Nenhum usuário isolado na rede.")
    else:
        for idx in indices_isolados:
            print(f"- {nomes[idx]}")

\end{minted}
\end{listing}


\begin{minted}[bgcolor = DarkBackground, linenos=true, breaklines]{Python}
def exibir_recomendacoes(nomes, matriz_A, matriz_A2, indice_alvo):
    """Lê a linha do alvo na matriz_A2 e sugere amigos em comum (que ainda não sejam amigos)."""
    nome_alvo = nomes[indice_alvo]
    print(f"--- SUGESTÕES DE AMIZADE PARA: {nome_alvo.upper()} ---")

    linha_do_alvo_A2 = matriz_A2[indice_alvo]
    linha_do_alvo_A = matriz_A[indice_alvo] # Precisamos ver as amizades originais
    tem_sugestao = False

    # O enumerate pega o índice (j) e o valor (amigos_em_comum) de uma vez só!
    for j, amigos_em_comum in enumerate(linha_do_alvo_A2):

        ja_sao_amigos = linha_do_alvo_A[j] > 0 # Verifica na matriz original

        # Só sugere se: tem amigos em comum, NÃO é a própria pessoa, e NÃO são amigos ainda
        if amigos_em_comum > 0 and j != indice_alvo and not ja_sao_amigos:
            print(f"Você pode se interessar por {nomes[j]} ({amigos_em_comum} pessoas que você segue também seguem esse perfil)")
            tem_sugestao = True

    if not tem_sugestao:
        print("Nenhuma sugestão no momento.")
\end{minted}

\begin{listing}[H]
	\begin{minted}[bgcolor = DarkBackground, linenos=true, breaklines]{Python}
# =============================================
# SCRIPT PRINCIPAL: APP SOCIAL (Test Harness)
# =============================================
import motor_matrizes as math_engine
import algoritmo_rede as rede
import interface_usuario as ui

# 1. Base de Dados
nomes = ["Ana", "Bruno", "Carlos", "Diana", "Eduardo"]

# Matriz de Conexões (A)
matriz_A = [
    [0, 1, 1, 0, 0],  # Ana segue Bruno e Carlos
    [1, 0, 1, 1, 0],  # Bruno segue Ana, Carlos e Diana
    [1, 0, 0, 1, 0],  # Carlos segue Ana e Diana
    [1, 0, 0, 0, 0],  # Diana segue Ana
    [0, 0, 0, 0, 0]   # Eduardo é isolado
]

print("=" * 40)
print("Beep Beep... INICIANDO ANÁLISE DA REDE SOCIAL ")
print("=" * 40 + "\n")
\end{minted}
\end{listing}


\begin{minted}[bgcolor = DarkBackground, linenos=true, breaklines]{Python}
try:
    # 2. Processamento (Lógica)
    indice_famoso = rede.encontrar_influenciador(matriz_A)

    # Proteção: só calcula seguidores se o aluno já programou a função e não retornou -1
    if indice_famoso != -1:
        seguidores_famoso = math_engine.somar_coluna(matriz_A, indice_famoso)
    else:
        seguidores_famoso = 0

    isolados = rede.encontrar_isolados(matriz_A)
    matriz_conexoes = rede.gerar_recomendacoes(matriz_A) # Isso deve retornar a matriz A^2

    # 3. Exibição (Interface)
    ui.exibir_perfil_usuario(nomes, indice_famoso, seguidores_famoso)
    print()
    ui.exibir_isolados(nomes, isolados)
    print()

    # Testando recomendações para a Ana (índice 0)
    if len(matriz_conexoes) > 0:
        # CORREÇÃO AQUI: Passando matriz_A e matriz_conexoes (A^2)
        ui.exibir_recomendacoes(nomes, matriz_A, matriz_conexoes, 0)
    else:
        print("A matriz de recomendações ainda está vazia. Implemente a multiplicação!")

except Exception as e:
    print(f"ERRO NA EXECUÇÃO: {e}")
    print("Verifique se você implementou as funções corretamente nos módulos!")
\end{minted}

\end{document}
